%% %&pdflatex
%% \RequirePackage[l2tabu, orthodox]{nag}

% article, report, book, ..., scrartcl, scrreprt, scrbook
\documentclass[9pt,a4paper,onecolumn,oneside]{extarticle}

%-- Language and font packages
\usepackage[utf8]{inputenc}    % Umlaute; default: ascii, latin1
\usepackage[T1]{fontenc}       % European languages support (256 Byte)
%\usepackage{fixltx2e}         % deprecated LaTeX2e improvements
\usepackage{lmodern}           % Latin modern fonts, not computer modern

\usepackage{pdflscape}
\usepackage{multicol}
%\setlength{\columnseprule}{0.4pt}
\setlength{\columnsep}{0.8cm}
%% \usepackage[margin=0.5in]{geometry}
\usepackage{amssymb,amsmath}

%\usepackage[light]{kpfonts}   % for "lighter" printing
%\usepackage[poorman]{fourier} 

%\usepackage[english,ngerman]{babel}  % 'australian' for European date

% Postscript fonts               % Times etc. obsolete !
%\usepackage{kmath,kerkis}       % Bookman
%\usepackage{fouriernc}          % New Century Schoolbook
%\usepackage{mathpazo}           % Palatino
%\usepackage{mathptmx}           % Times Roman
%\usepackage{ccfonts,eulervm}    % Concrete and Euler fonts
% or:
%\renewcommand{\rmdefault}{phv}
%\renewcommand{\sfdefault}{phv}  % Helvetica

\usepackage{microtype}          % improves spacing between words and letters
%\usepackage{xspace}             % \newcommand{\eg}{e.\,g.\xspace}
%\usepackage{ulem}               % \u[u]line \sout \xout, but: \emph redefined

%-- Page layout and headings
\usepackage[margin=1.5cm]{geometry}
%\usepackage[top=2.5cm, bottom=3.0cm, left=4.0cm, right=4.0cm]{geometry}
%\usepackage{fullpage}  % or: \usepackage{a4wide}
%\usepackage[a4paper]{geometry}
%\usepackage{layouts}  % \showpage

%\usepackage[compact,small]{titlesec}    %\itemsep 0pt
%\usepackage{sectsty}
%\sectionfont{
%  \sectionrule{0pt}{0pt}{-5pt}{0.8pt}  % underline 0.8pt with distance 5pt
%}

\usepackage{paralist}   % \begin{compactitem}...\item...\end{compactitem}
%\usepackage{mdwlist}   % \begin{enumerate*,itemize*}...\item...\end{...}
%\usepackage{multicol}  % \begin{multicols}{2}...\end{multicols}

%-- Footer and header lines
\pagestyle{empty}  % empty, plain, headings, myheadings
%\markright{Heading \hrulefill\ }
%\markboth{Hans W. Borchers \hrulefill\ }{ \hrulefill\ Hans W. Borchers}
%\usepackage{fancyhdr,lastpage}
%\setlength{\headheight}{15pt}
%\pagestyle{fancy}
%\fancyhead{} \fancyfoot{}  % or \fancyhf{}: remove everything
%\lhead{\textsc{Project Title}} \chead{\textsl{Draft}} \rhead{}
%\lfoot{\small Hans W. Borchers \copyright{} 2010} \cfoot{}
%\rfoot{Page \thepage\ of \pageref{LastPage}}
%\renewcommand\headrulewidth{0pt}  % 1.2pt
%\renewcommand\footrulewidth{0pt}  % 0.4pt

%-- Miscellaneous, hyphenation, etc.
% \setlength{\parindent}{0pt} \setlength{\parskip}{6pt plus3pt minus2pt}
\usepackage[parfill]{parskip}            
\usepackage{todonotes}
\hyphenation{FORTRAN Hy-phen-a-tion}  % or: mbox{FORTRAN}; use \- in text
%\usepackage{makeidx}

%-- Including graphics
\usepackage{graphicx}
%\DeclareGraphicsExtensions{.pdf,.png,.jpg}
%\graphicspath{{.}}

%-- 'tabularx' and other environments
\usepackage{tabularx}
%\usepackage{booktabs}  % professionel
%\usepackage{supertabular}
%\usepackage{longtables}

%-- Citation styles
%\usepackage[backend=biber,sortlocale=de,natbib=true]{biblatex}
%\usepackage{natbib}  % and use citet{},...
%\bibliographystyle{plainnat}
%\bibpunct{(}{)}{;}{a}{,}{,}

%-- Math packages
%\usepackage{mathtools} %\usepackage{amsmath,amssymb,amsfonts}

%-- Algorithms and listings
\usepackage{alltt}              % 'alltt' environment
\usepackage{relsize,fancyvrb}   % 'Verbatim' environment
\DefineShortVerb{\|}            % \UndefineShortVerb{\|}
\usepackage{listings}

%-- Packages for including symbols
%\usepackage[official]{eurosym}  % or: [gen]; \euro
\usepackage{textcomp}   % \texteuro and \oldstylenums{01234567890}
\usepackage{textgreek}  % \textDelta\textbeta
\usepackage{pifont}     % \ding{32..126, 161..254}

%-- Utilitiy packages
\usepackage{color}
\definecolor{darkred}{rgb}{0.5,0.0,0.0}

%\usepackage{calc}
%\usepackage{csquotes}  % context sensitive quotation
%\usepackage{siunitx}   % \num, \si, ex.: \SI{10}{\hertz}, package 'SIunits'
\usepackage{url}

%-- Last packages to load
%\usepackage[colorlinks=false, pdfborder={0 0 0}]{hyperref}
%\hypersetup{colorlinks=true, linkcolor=blue, filecolor=blue, pagecolor=blue, urlcolor=blue}
%\usepackage{cleveref}  % \cref instead of \ref or \eqref

%-- Utility settings
\tolerance 1414
\hbadness 1414
\emergencystretch 1.5em
\clubpenalty = 10000    % prevent orphans and widows
\widowpenalty=10000
\vfuzz \hfuzz           % \hfuzz 0.3pt
\raggedbottom           % \sloppy --- is l2tabu
\frenchspacing

%-- New command definitions    
%\newcommand{\zB}{\mbox{z.\,B.}\xspace}    
%\newcommand{\tag}[1]{\item [\em #1]\mbox{}  \\}
%\newcommand{\tick}[1]{\item {\em #1}}
%\newfont{\dunhill}{cmdunh10 scaled 1000}

%-- Changed commands and settings
%\renewcommand{\abstractname}{Executive Summary}
\setcounter{secnumdepth}{2}
\setcounter{tocdepth}{2}

%-- --------------------------------------------------------------------
\begin{document}
\typeout{Copyright (c) 2000--2012, Hans W. Borchers}%
\begin{landscape}
%% Put document material here. Note title setup

\includegraphics[height=0.4in]{Rlogo.png} 
\hspace{5mm}
\large{Base R Optimization Cheat Sheet --- Version 2020-09-28 \hfill JC Nash and HW Borchers (R-optimization at mailbox.org)}
%\large{Base R Optimization Cheat Sheet \hspace*{9mm}\small{J C Nash  nashjc and H W Borchers -- 2020/09/27}}

\begin{multicols}{2}

\section*{\color{darkred} Purpose and Formulation}

This cheatsheet is a quick guide to optimization tools present in the base 
distribution of the R software for statistical computing.

\textbf{General problem}:

Find a vector $\mathbf{x}$ of $n$ parameters that minimizes a given function 
$f(\mathbf{x})$ in a certain domain, i.e. subject to inequality constraints 
$g(\mathbf{x}) \le 0$, and sometimes also equality constraints $h(\mathbf{x}) = 0$.

Numerical optimization routines will in general only find \textbf{local} minima 
that are minimal only in their neighborhood. There may exist several or even 
many minima in a given domain.

\begin{itemize}
%\tightlist
\item
  Most problems are special cases. Base R mainly treats 
  problems that are unconstrained minimizations of a few parameters.
\item
  There are often several ways to formulate a problem for different software  
  tools, and some of these may be much more efficient than others.
\item
  Bounds constraints \textbf{lb} $\le$ \textbf{x} $\le$
  \textbf{ub} are a special and easier to handle case of constraints.
\item
  To \textit{maximize} a function f(\textbf{x}), minimize -f(\textbf{x}) 
\item
  Most tools require suitable starting parameters \textbf{x0}.
  This is particularly important if there are multiple minima.
  The minimum found is not guaranteed to be the nearest one to |x0|.
\item
  The lowest of all minima is known as the \textbf{global} minimum.\\
  Unfortunately, Base R has no tools for global optimization.
\item
  Iterative algorithms \textbf{converge}, but programs
  \textbf{terminate} (convergence, computational failure, limit on
  effort).
\item
  The tools in Base R are intended for smooth functions. 
  Optimizing non-smooth functions usually requires special approaches, though
  \texttt{optim()} with method \texttt{Nelder-Mead} may sometimes succeed.
\end{itemize}

For many more optimization tools see the Task View for Optimization.

REMARK: In R syntax below, square brackets indicate an optional element of a command.
\texttt{...} indicates additional parameters of the function to be minimized.
%\vfill
%\columnbreak


\section*{\color{darkred}One-dimensional optimization (x a scalar)}

\begin{alltt}
 result <- optimize(f, interval=c(lower, upper), ...,
                    tol=.Machine\$double\^{}0.25)
\end{alltt}

\begin{itemize}
\item
  \texttt{f} is a function of just one parameter, and \texttt{interval}
  gives the end points of the domain in which a minimum is to be found.
\item
  See the documentation for alternate syntax, including \texttt{maximum=TRUE}.
  The \texttt{tol} above is the default value (approx. 0.0001). This default
  tolerance is not small enough for many applications.
\item
  \texttt{result} is a list with elements \texttt{minimum} and
  \texttt{objective}.
\item
  If there are multiple minima, you only get one of them.
  If there are no minima -- you still get a result! CAUTION.
\item
  In one dimension it is computationally feasible to identify
  \textit{all} minima (or maxima) and thus also the global one(s).
\end{itemize}

There are many more one-dimensional minimizers in packages on CRAN.


\section*{\color{darkred} Unconstrained minimization}

Function \texttt{f(x [, ...])} computes the function to be minimized at 
\textbf{x}; choose a starting parameter \textbf{x0}. Optional, but
useful: a vector valued function \texttt{grad()} that computes the gradient
of \texttt{f()} at \textbf{x}.

|optim| is the main optimizer function in Base R. There are two other
minimizers that may be sometimes useful, \verb|nlm| and \verb|nlminb|.

\begin{alltt}
 result <- optim(x0, f [, grad] [, method = "Method"],
                 control = list())
\end{alltt}
where \texttt{Method} is one of 'Nelder-Mead', 'BFGS', 'L-BFGS-B', 'CG', or 'SANN'.
The default is 'Nelder-Mead', a gradient-free -- but slow -- optimization solver.

NOTE: \texttt{CG} (JN is author!) and \texttt{SANN} are NOT recommended.

The most important control option is the \textit{tolerance}, called 
\texttt{reltol} for methods 'Nelder-Mead' and 'BFGS', and \texttt{factr} 
for method 'L-BFGS-B'. Its default value is about |1e-08|.

\vfill
\columnbreak


|nlm| carries out a Newton-type algorithm and returns a list with |minimum| 
for the function value and |estimate| for the solution vector.

\begin{alltt}
 result <- nlm(f, x0)
\end{alltt}

\begin{itemize}
 \item
A gradient function is not used. Look at the help page to see, how the
gradient is to be defined and used in \verb|nlm|. Also note the order of
the call for \texttt{nlm} has the function first, then the starting parameters.
\end{itemize}

|nlminb| uses the PORT routines, a FORTRAN implementation of quasi-Newton BFGS.
It has been included with Base R for historical reasons.

\begin{alltt}
 result <- nlminb(x0, f [, grad])
\end{alltt}

It returns |par| and |objective| as list (with some more information).

REMARKS:

\begin{itemize}
\item
  To minimize functions of many parameters, use \texttt{optim()} with method 
  \texttt{L-BFGS-B} or find efficient solvers on the Task View for Optimization.
\item
  When a gradient is needed but not specified, methods use numerical
  approximations (finite differences). This generally gives slightly less
  satisfactory solutions and may require more computational effort.
\end{itemize}



\section*{\color{darkred} Bounds constrained optimization}

From \texttt{optim()}, only method \texttt{L-BFGS-B} can deal with
bounds. \texttt{nlminb()} also handles bounds. Assume variables
\texttt{lb} and \texttt{ub} have either single values or else as many 
values as \textbf{x} to define lower and upper bounds respectively.

\begin{alltt}
 result <- optim(x0, f [,grad], method = "L-BFGS-B",
                 lower = lb, upper = ub)
\end{alltt}

\begin{alltt}
 result <- nlminb(x0, f [, grad], lower = lb, upper = ub)
\end{alltt}

\begin{itemize}
\item
  From a practical viewpoint, methods often fail to find appropriate
  solutions when lower and upper bounds are specified close together.
\item
  Bounds constrained optimization without the specification of a
  gradient function may lead to computational failure when it tries 
  to evaluate the function outside the bounds.
\end{itemize}


\section*{\color{darkred} Linearly constrained optimization}

The function \texttt{constrOptim} will approximate \textit{linearly} 
constrained optimization problems, where there are k linear constraints 
such that \textbf{A} \%*\% \textbf{x} $\ge$ \textbf{b}
with \textbf{A} a k-by-n matrix, n the number of parameters and
\textbf{b} a length k vector of constants.

\begin{alltt}
 result <- constrOptim(x0, f [, grad], ui=A, ci=b [, method])
\end{alltt}

\begin{itemize}
\item
  \texttt{constrOptim} uses an adaptive barrier method in conjuction
  with \texttt{optim()}, so it has the same strengths and weaknesses as
  \texttt{optim()}.
\item
  Be careful that the dimensions of \texttt{A} and \texttt{b} are set
  up correctly.
\item
  The default method is 'BFGS', 'L-BFGS-B is likely more efficient.
\end{itemize}

NOTE: In Base R there are no optimizers for nonlinear inequality constraints
or even for equality constraints. See the Optimization Task View for CRAN 
packages solving such problems.



\section*{\color{darkred} Nonlinear least squares}

If \texttt{f()} is a sum of squared terms, then special methods can be
used. Base R has the \texttt{nls()} function for such problems. It has
many options. Here we will just give an example of the most common task,
which is modeling.

Suppose we have a data frame (or list) \texttt{Data} with columns |x| and
(the dependent) |y|. The model shall be a formula
\( y \sim a + b \cdot e^{c x} \)
with parameters \( a, b, c \) to be determined from the data. The call to 
\texttt{nls()} looks likely

\begin{alltt}
 result <- nls(y ~ a + b*exp(c*x), data = Data,
               start = c(a=1.0, b=2.0, c=0.05)
               [, lower = c(0,0,0), upper = c(5.0,5.0,1.0)]
               [, control = list()] [, algorithm = "port"])
\end{alltt}

If there are bounds constraints, the algorithm must be |port|, else it 
can be |plinear| or |NULL| (the default), a Gauss-Newton type of approach.
|start| will be a named vector with starting values for all unknown 
parameters.

\begin{itemize}
%\tightlist
\item
  \texttt{nls()} is a powerful tool in the hands of an experienced user,
  but often leads to ``singular gradient'' errors, that means it has been 
  unable to find a solutions. Packages \texttt{nlsr} and \texttt{minpack.lm} 
  should be considered.
\end{itemize}


\section*{\color{darkred} See also}\label{See also}

CRAN Task View on Optimization: \\
\hspace*{9mm}\url{https://cran.r-project.org/web/views/Optimization.html} 


\end{multicols}

\end{landscape}
\end{document}
